% Source: Wald GR, physical PDF pp. 79--87
%         (printed pp. 66--74).
% Coverage: Section 4.3, General Relativity; equations (4.3.1)--(4.3.23).
% Figures/tables/footnotes: footnotes 5--6; no figures or tables.
% Status: source-reviewed and compile-clean.
% Uncertainties: none.

\subsection{General Relativity}\label{sec:4.3}

Maxwell's theory is a remarkably successful theory of electricity, magnetism, and
light which is beautifully incorporated into the framework of special relativity.
Therefore, one might expect that the next logical step would have been to develop
a new theory of the other classical force, gravitation, which would generalize
Newton's theory and make it compatible with special relativity in the same way that
Maxwell's theory generalized Coulomb's electrostatics. However, Einstein chose an
entirely different path and instead developed general relativity, a new theory of
spacetime structure and gravitation. As already mentioned in the introductory
chapter, the equivalence principle and Mach's principle provided the primary
motivation for formulating a new theory.

To see the relevance of the equivalence principle---that all bodies fall the same
way in a gravitational field---to developing a new viewpoint on gravitation,
consider how one measures the electromagnetic field in special relativity. The
first step is to set up \textquotedblleft background
observers\textquotedblright{} who are not subject to electromagnetic forces (i.e.,
they are electrically neutral, have no magnetic dipole moment, etc.) or any other
forces. These observers are called inertial and satisfy the geodesic equation of
motion~\eqref{eq:4.2.6}. The next step is to release a charged test body. The
world line of this body will satisfy equation~\eqref{eq:4.2.26}, and by observing
the deviation from inertial motion (for sufficiently many test bodies) we can
determine \(F_{ab}\).

If we apply this procedure to gravitation, we are immediately faced with a serious
problem: By the equivalence principle, we have no way of
\textquotedblleft insulating\textquotedblright{} an observer or body from the
gravitational force, so we have no simple, direct physical procedure for
constructing inertial observers in the sense used for electromagnetism. Any
observer will move in exactly the same way as a test body, so we have no natural
\textquotedblleft background motion\textquotedblright{} to compare with the test
body. Thus we have no simple, direct way of measuring the gravitational force
field. It is, of course, possible that complicated procedures for accomplishing
these constructions and measurements may exist. If special relativity were
correct, one could construct inertial observers by spacetime measurements, i.e.,
the (flat) spacetime metric could be measured using clocks and metersticks, and its
geodesics determined. Inertial observers would need to be equipped with rocket
engines, but aside from that the gravitational force field could be determined in
the same way as for electromagnetism. The equivalence principle would then be
viewed as a peculiar quirk of the gravitational force law, just as it is viewed in
standard treatments of Newtonian theory.

The basic framework of the theory of general relativity arises from considering the
opposite possibility: that we cannot in principle---even by complicated
procedures---construct inertial observers in the sense of special relativity and
measure the gravitational force. This is accomplished by the following bold
hypothesis: \emph{The spacetime metric is not flat, as was assumed in special
relativity. The world lines of freely falling bodies in a gravitational field are
simply the geodesics of the (curved) spacetime metric.} In this way, the
\textquotedblleft background observers\textquotedblright{} (geodesics of the
spacetime metric) automatically coincide with what was previously viewed as
motion in a gravitational force field. As a result we have no meaningful way of
describing gravity as a force field; rather, we are forced to view gravity as an
aspect of spacetime structure. Although absolute gravitational force has no
meaning, the relative gravitational force (i.e., tidal force) between two nearby
points still has meaning and can be measured by observing the relative acceleration
of two freely falling bodies. This relative acceleration is directly related to the
curvature of spacetime by the geodesic deviation equation (3.3.18).

How does this viewpoint of general relativity that there is no such thing as
gravitational force square with the well known \textquotedblleft
fact\textquotedblright{} that there is a gravitational force field at the surface of
the Earth of \(980\,\mathrm{cm\,s^{-2}}\)? Recall that in the standard Newtonian
viewpoint this gravitational force on an object placed on the Earth's surface is
balanced by the force the surface exerts, leaving the body in equilibrium, i.e.,
\textquotedblleft at rest\textquotedblright. In the viewpoint of general relativity,
the only force acting on the body is the force of the surface of the Earth. On
account of this force, the body accelerates (i.e., deviates from geodesic motion) at
the rate of \(980\,\mathrm{cm\,s^{-2}}\). Nevertheless, it remains in a stationary
state, because in the curved spacetime geometry in the vicinity of the Earth, the
orbits of time translation symmetry differ from the geodesics of the metric. We
could use the time translation symmetry of this example to define a preferred set of
background observers. We then could define the gravitational force field of the
Earth to be minus the acceleration a body must undergo in order to remain
stationary. Thus, in this case a well defined meaning can be assigned to gravity as a
force field. However, in the absence of time translation symmetry---e.g., in a case
where there are several massive bodies in relative motion---there exists no natural
set of curves whose comparison with geodesics could be used to define
gravitational force. (This remark also applies to an observer in a nonstationary
laboratory near the Earth who would not be aware of the time translation
symmetry.) Thus, although we may meaningfully speak of the gravitational force
field of the Earth, this is a very special notion, applicable only to situations with
time translation symmetry. In general situations, there are no preferred
background observers, and only the tidal force---the relative acceleration of nearby
geodesics---is well defined.

Thus, in general relativity, we do not assert that spacetime is the manifold
\(\mathbb R^4\) with a flat metric \(\eta_{ab}\) defined on it. This is a
possibility corresponding to no tidal forces, i.e., no gravitational field, but it is
not the only possibility. The framework of general relativity permits the Lorentz
metric, \(g_{ab}\), of spacetime to be curved. Indeed, it asserts that spacetime
must be curved in all situations where, physically, a gravitational field is present.
Since we are allowing curved geometries, it is also much more natural to allow
spacetime to have a manifold structure \(M\) other than \(\mathbb R^4\). Hence,
in general relativity we place no \emph{a priori} restriction on the spacetime
manifold. The final crucial feature of general relativity is Einstein's equation
which relates the spacetime geometry to the matter distribution. We will discuss
this equation below, but first we discuss the nature of the laws of physics under the
new framework of spacetime structure given by general relativity:
\emph{Spacetime is a manifold \(M\) on which is defined a Lorentz metric
\(g_{ab}\).}

The laws of physics in general relativity are governed by two basic principles:
(1) the principle of general covariance---already discussed in the previous two
sections---which states that the metric, \(g_{ab}\), and quantities derivable from
it are the only spacetime quantities that can appear in the equations of physics;
(2) the requirement that equations must reduce to the equations satisfied in
special relativity in the case where \(g_{ab}\) is flat. As previously mentioned,
the first principle is imprecise because the term \textquotedblleft spacetime
quantity\textquotedblright{} is not well defined. As will be illustrated below,
these two principles alone do not uniquely determine the laws of physics in general
relativity. However, together with simplicity and aesthetics, they serve as guides
which, in many cases, lead directly to natural candidates for physical laws.

Since the basic framework of general relativity modifies that of special relativity
only in that it allows the manifold to differ from \(\mathbb R^4\) and the metric
to be nonflat, we may continue to represent physical quantities by the same type of
tensor fields as in special relativity. Thus, in general relativity, particle motion
continues to be represented by a timelike (in the metric \(g_{ab}\)) curve; perfect
fluids are still described in terms of a 4-velocity \(u^a\), a density \(\rho\),
and a pressure \(P\); the electromagnetic field is represented by an
antisymmetric tensor \(F_{ab}\). Only the equations satisfied by these fields
need to be amended. The above two principles suggest the following simple rule:
In the equations holding in special relativity, replace everywhere the metric
\(\eta_{ab}\) of special relativity by \(g_{ab}\) and correspondingly replace
the derivative operator \(\partial_a\) associated with \(\eta_{ab}\) by the
derivative operator \(\nabla_a\) associated with \(g_{ab}\). This rule for, in
effect, coupling particles and fields to gravity is closely analogous to the
\textquotedblleft minimal coupling\textquotedblright{} rule
\(p_a\longrightarrow p_a-eA_a\) for coupling to electromagnetism. However,
as we shall illustrate below, this rule is not entirely free of ambiguity.

Thus, in general relativity, we again define the 4-velocity, \(u^a\), of a
particle to be the unit tangent (as measured by \(g_{ab}\)) to its world line. A
free particle satisfies the geodesic equation of motion,
\begin{equation}
  u^a\nabla_au^b=0 ,
  \tag{4.3.1}\label{eq:4.3.1}
\end{equation}
where \(\nabla_a\) is the derivative operator associated with \(g_{ab}\). If
the acceleration\footnote{The (absolute) acceleration \(a^b\) should be clearly
distinguished from the relative acceleration of geodesics discussed in section 3.3.}
\(a^b=u^a\nabla_au^b\) of the particle is nonvanishing, we say that a force
\(f^b=ma^b\) acts on the particle, where \(m\) is its (rest) mass. For
example, if the particle has (rest) mass \(m\) and charge \(q\), and is placed
in an electromagnetic field \(F_{ab}\), it satisfies the Lorentz force equation,
\begin{equation}
  u^a\nabla_au^b=\frac{q}{m}F^b{}_cu^c ,
  \tag{4.3.2}\label{eq:4.3.2}
\end{equation}
where indices are raised and lowered by \(g_{ab}\), i.e.,
\(F^b{}_c=g^{bd}F_{dc}\). (Again, we emphasize that there is no natural flat
metric defined on spacetime and, by the principle of general covariance, only
\(g_{ab}\) can enter the equations.) The 4-momentum of the particle is defined
by
\begin{equation}
  p^a=mu^a .
  \tag{4.3.3}\label{eq:4.3.3}
\end{equation}
The energy of the particle as determined by an observer who is present at the
event on the particle's world line at which the energy is measured is again
\begin{equation}
  E=-p_av^a ,
  \tag{4.3.4}\label{eq:4.3.4}
\end{equation}
where \(v^a\) is the 4-velocity of the observer. However, there is one important
difference here: Because spacetime is curved, there is no well defined notion of
vectors at different points being parallel; parallel transport is curve dependent.
Thus, there is no natural \textquotedblleft global family\textquotedblright{} of
inertial observers, and a given observer cannot, in general, define the energy of a
distant particle.

In general relativity, continuous matter distributions and fields again are described
by a stress-energy tensor \(T_{ab}\). The stress tensor of a perfect fluid is given
by
\begin{equation}
  T_{ab}=\rho u_au_b+P(g_{ab}+u_au_b)
  \tag{4.3.5}\label{eq:4.3.5}
\end{equation}
and it satisfies the equations of motion
\begin{equation}
  \nabla^aT_{ab}=0 ,
  \tag{4.3.6}\label{eq:4.3.6}
\end{equation}
which yield
\begin{align}
  u^a\nabla_a\rho+(\rho+P)\nabla^au_a&=0 ,
    \tag{4.3.7}\label{eq:4.3.7}\\
  (P+\rho)u^a\nabla_au_b+(g_{ab}+u_au_b)\nabla^aP&=0 .
    \tag{4.3.8}\label{eq:4.3.8}
\end{align}
However, the interpretation of equation~\eqref{eq:4.3.6} is altered now. A
family of observers is represented by a unit timelike vector field \(v^a\). If
one could find such a vector field which is covariantly constant, i.e.,
\(\nabla_av_b=0\)---or for which merely \(\nabla_{(a}v_{b)}=0\)---then we
would have \(\nabla^a(T_{ab}v^b)=0\). Applying the curved spacetime version
of Gauss's law (see appendix B), we again would obtain strict conservation of
energy in the form~\eqref{eq:4.2.18} for the energy-momentum four-vector
\(J_a=-T_{ab}v^b\) measured by the observers represented by \(v^b\).
However, in curved spacetime in general one no longer can find a \(v^a\)
satisfying \(v^av_a=-1\) and \(\nabla_{(a}v_{b)}=0\). (Indeed, the
equation \(\nabla_{(a}v_{b)}=0\) is Killing's equation and holds if and only
if \(v^a\) generates a one-parameter group of isometries [see appendix C].)
Thus the argument fails that equation~\eqref{eq:4.3.6} implies strict energy
conservation. Physically, this makes sense because the gravitational tidal forces
can do work on the fluid and may increase or decrease its locally measured
energy.\footnote{One might hope to recover an energy conservation law by including
the stress-energy of the gravitational field as can be done in Newtonian theory.
However, in general relativity there exists no meaningful local expression for
gravitational stress-energy and thus there is no meaningful local conservation law
which leads to a statement of energy conservation. Nevertheless, as will be
discussed in chapter 11, a conserved total energy of an isolated system can be
defined, even though there is no local expression for energy density.} However, if
one considers a spacetime region of dimension small compared with radii of
curvature, then, physically, the tidal forces can do little work and the energy of
the fluid should be approximately conserved. But over this small spacetime region
it is possible to find vector fields with \(\nabla_bv^a\approx0\), and thus
equation~\eqref{eq:4.3.6} does yield approximate conservation of energy as
measured by these observers. Thus, equation~\eqref{eq:4.3.6} may be
interpreted as a local conservation of material energy over small regions of
spacetime. On account of this interpretation, we expect
equation~\eqref{eq:4.3.6} to hold for all matter and fields, not just for perfect
fluids.

The most natural generalization of the equation satisfied by a Klein-Gordon scalar
field to curved spacetime is given by our \textquotedblleft minimal
substitution\textquotedblright{} rule
\(\eta_{ab}\longrightarrow g_{ab}\), \(\partial_a\longrightarrow\nabla_a\),
\begin{equation}
  \nabla^a\nabla_a\phi-m^2\phi=0 .
  \tag{4.3.9}\label{eq:4.3.9}
\end{equation}
The stress tensor of the field is
\begin{equation}
  T_{ab}=\nabla_a\phi\,\nabla_b\phi
  -\frac12g_{ab}\left(\nabla_c\phi\,\nabla^c\phi+m^2\phi^2\right)
  \tag{4.3.10}\label{eq:4.3.10}
\end{equation}
and satisfies \(\nabla^aT_{ab}=0\). We should point out, however, that there
are many other possible generalizations of equation~\eqref{eq:4.2.19} which
are consistent with the two basic principles stated above. For example, the
equation
\begin{equation}
  \nabla^a\nabla_a\phi-m^2\phi-\alpha R\phi=0 ,
  \tag{4.3.11}\label{eq:4.3.11}
\end{equation}
where \(\alpha\) is a constant, is such a generalization. Indeed,
equation~\eqref{eq:4.3.11} with \(\alpha=1/6\) arises naturally on account of
its conformal invariance properties (see appendix D).

Maxwell's equations in curved spacetime become
\begin{align}
  \nabla^aF_{ab}&=-4\pi j_b ,
    \tag{4.3.12}\label{eq:4.3.12}\\
  \nabla_{[a}F_{bc]}&=0 .
    \tag{4.3.13}\label{eq:4.3.13}
\end{align}
The electromagnetic stress tensor again is given by
equation~\eqref{eq:4.2.27} with \(g_{ab}\) replacing \(\eta_{ab}\),
\begin{equation}
  T_{ab}=\frac{1}{4\pi}\left\{
  F_{ac}F_b{}^c-\frac14g_{ab}F_{de}F^{de}\right\} .
  \tag{4.3.14}\label{eq:4.3.14}
\end{equation}
Again, equation~\eqref{eq:4.3.13} allows us to introduce a vector potential
\(A_a\) (at least locally). However, Maxwell's equations for \(A_a\) in the
Lorenz gauge contains an explicit curvature term resulting from the commutation
of derivatives in the derivation of equation~\eqref{eq:4.2.32}; we find
\begin{equation}
  \nabla^a\nabla_aA_b-R^d{}_bA_d=-4\pi j_b .
  \tag{4.3.15}\label{eq:4.3.15}
\end{equation}
This illustrates an important deficiency of our minimal substitution rule. Had we
minimally substituted in Maxwell's equations in the form~\eqref{eq:4.2.32},
we would have been led to equation~\eqref{eq:4.3.15} without the Ricci tensor
term. In this instance, we can decide in favor of equation~\eqref{eq:4.3.15}
over the alternative equation without the \(R^d{}_bA_d\) term because
equation~\eqref{eq:4.3.15} implies current conservation,
\(\nabla_aj^a=0\) (problem 1), while the alternative equation conflicts with it.
However, this example shows that \textquotedblleft minimal
substitution\textquotedblright{} by itself is not a unique prescription.

In situations where the spacetime scale of variation of the electromagnetic field
is much smaller than that of the curvature, one would expect to have solutions of
Maxwell's equations of the form of a wave oscillating with nearly constant
amplitude, i.e., solutions of the form
\begin{equation}
  A_a=C_ae^{iS} ,
  \tag{4.3.16}\label{eq:4.3.16}
\end{equation}
where derivatives of \(C_a\) are \textquotedblleft small\textquotedblright.
Substituting equation~\eqref{eq:4.3.16} into equation~\eqref{eq:4.3.15}
with \(j_b=0\) and neglecting the \textquotedblleft small\textquotedblright{}
term \(\nabla_b\nabla^bC_a\) as well as the Ricci tensor term yields the
condition
\begin{equation}
  \nabla_aS\,\nabla^aS=0 ;
  \tag{4.3.17}\label{eq:4.3.17}
\end{equation}
i.e., we find again that the surfaces of constant phase are null, and thus (by the
same argument as given above for flat spacetime) \(k_a=\nabla_aS\) is tangent
to null geodesics. This suggests that, in this approximation (known as the
\emph{geometrical optics approximation}), light travels on null geodesics, a
suggestion which can be confirmed by studies of the Green's function.

We now have described how general relativity treats gravitation in terms of curved
spacetime geometry and we have illustrated the nature of the laws of physics in this
new framework of spacetime structure. The remaining ingredient of general
relativity is the equation satisfied by the spacetime metric. It is here that Mach's
principle comes into play. Rather than prescribe the spacetime geometry in advance,
general relativity asserts that the spacetime geometry is influenced by the matter
distribution in the universe, in accordance with some of Mach's ideas (see section
1.4). In this way the spacetime metric now becomes not only a background arena
on which the laws of physics are staged but also a dynamical variable which
responds to the matter content of spacetime, as must be the case if the spacetime
geometry is to describe gravity.

What equation describes the relation between spacetime geometry and the matter
distribution? An important clue is provided by the comparison of the description of
tidal force in Newtonian gravity and general relativity. In the Newtonian theory, the
gravitational field may be represented by a potential, \(\phi\), and the tidal
acceleration of two nearby particles is given by
\(-(\mathbf{x}\mathbin{\cdot}\boldsymbol{\nabla})\boldsymbol{\nabla}\phi\), where \(\mathbf{x}\) is
the separation vector of the particles. On the other hand, in general relativity,
from equation~(3.3.18) the tidal acceleration of two nearby particles is given by
\(-R_{cbd}{}^av^cx^bv^d\), where \(v^a\) is the 4-velocity of the particles
and \(x^a\) is the deviation vector. This suggests that we make the
correspondence,
\begin{equation}
  R_{cbd}{}^av^cv^d\longleftrightarrow\partial_b\partial^a\phi .
  \tag{4.3.18}\label{eq:4.3.18}
\end{equation}
However, Poisson's equation tells us that
\begin{equation}
  \boldsymbol{\nabla}^2\phi=4\pi\rho ,
  \tag{4.3.19}\label{eq:4.3.19}
\end{equation}
where \(\rho\) is the mass (i.e., energy) density of matter and we remind the
reader that we use units where \(G=c=1\) here and throughout the text.
Furthermore, as we have discussed above, in special and general relativity the
energy properties of matter are described by a stress-energy tensor \(T_{ab}\),
and we have the correspondence
\begin{equation}
  T_{ab}v^av^b\longleftrightarrow\rho ,
  \tag{4.3.20}\label{eq:4.3.20}
\end{equation}
where \(v^a\) is the 4-velocity of the observer.

The correspondences~\eqref{eq:4.3.18} and~\eqref{eq:4.3.20} together with
equation~\eqref{eq:4.3.19} suggest that we have
\(R_{cad}{}^av^cv^d=4\pi T_{cd}v^cv^d\), which suggests the field equation
\(R_{cd}=4\pi T_{cd}\). Indeed, this equation was originally postulated by
Einstein. However, it has a serious defect. As discussed above, the stress tensor
satisfies \(\nabla^cT_{cd}=0\). On the other hand, the contracted Bianchi
identity, equation~(3.2.31), tells us that
\(\nabla^c(R_{cd}-\frac12g_{cd}R)=0\). Hence equality of \(R_{cd}\) and
\(4\pi T_{cd}\) would imply \(\nabla_dR=0\), i.e., that \(R\), and hence
\(T=T^a{}_a\), is constant throughout the universe. This is a highly unphysical
restriction on the matter distribution, and it forces us to reject this equation, as
Einstein quickly realized (Einstein 1915b).

However, this difficulty also suggests its resolution. If instead we consider the
equation
\begin{equation}
  G_{ab}\equiv R_{ab}-\frac12Rg_{ab}=8\pi T_{ab} ,
  \tag{4.3.21}\label{eq:4.3.21}
\end{equation}
then there is no longer a conflict between the Bianchi identity and local
conservation of energy; indeed, the Bianchi identity implies local energy
conservation if equation~\eqref{eq:4.3.21} holds. Furthermore, the
correspondences which motivated the previous equation are not destroyed. Taking
the trace of equation~\eqref{eq:4.3.21}, we find
\begin{equation}
  R=-8\pi T ,
  \tag{4.3.22}\label{eq:4.3.22}
\end{equation}
and thus,
\begin{equation}
  R_{ab}=8\pi\left(T_{ab}-\frac12g_{ab}T\right) .
  \tag{4.3.23}\label{eq:4.3.23}
\end{equation}
In situations where Newtonian theory should be applicable, the energy of matter
as measured by an observer who is roughly \textquotedblleft at
rest\textquotedblright{} with respect to the masses will be much greater than the
material stresses (in units where \(c=1\)), so we have
\(T\approx-\rho=-T_{ab}v^av^b\). Thus, in this case,
equation~\eqref{eq:4.3.23} still leads to
\(R_{ab}v^av^b\approx4\pi T_{ab}v^av^b\).

Equation~\eqref{eq:4.3.21} is the desired field equation of general relativity. It
was written down by Einstein in 1915 and is known as \emph{Einstein's equation}.
The entire content of general relativity may be summarized as follows:
\emph{Spacetime is a manifold \(M\) on which there is defined a Lorentz metric
\(g_{ab}\). The curvature of \(g_{ab}\) is related to the matter distribution
in spacetime by Einstein's equation~\eqref{eq:4.3.21}.}

Most of the rest of this book is devoted to studying the solutions of Einstein's
equation and their physical properties. However, before concluding this section,
we make three brief remarks concerning the nature of this equation. The first
remark concerns its mathematical character. If we choose a coordinate system and
express the coordinate basis components \(R_{\mu\nu}\) in terms of
\(g_{\mu\nu}\), we see from section 3.4a that \(R_{\mu\nu}\) depends on
derivatives of \(g_{\mu\nu}\) up to second order, and is highly nonlinear in
\(g_{\mu\nu}\) (although it is linear in the second derivatives of
\(g_{\mu\nu}\)). Thus, Einstein's equation is equivalent to a coupled system of
nonlinear second order partial differential equations for the metric components
\(g_{\mu\nu}\). For a metric of Lorentz signature, these equations have a
hyperbolic (i.e., wave equation) character (see chapter 10). That we have the
correct number of equations and unknowns to permit a good initial value formulation
will be shown in chapter 10. Some methods for solving Einstein's equation will be
discussed in chapter 7.

The second remark concerns how one should view Einstein's equation. In one sense,
Einstein's equation~\eqref{eq:4.3.21} is analogous to Maxwell's
equation~\eqref{eq:4.2.32} with the stress tensor \(T_{ab}\) serving as the
source of the gravitational field in much the same way as the current \(j_a\)
serves as a source of the electromagnetic field. However, there is an important
difference. It makes sense to solve Maxwell's equation by specifying \(j_a\)
first, and then finding \(A_a\). One could try to solve Einstein's equation by
specifying \(T_{ab}\) first and then finding \(g_{ab}\). However, this does not
make much sense because until \(g_{ab}\) is known, we do not know how to
physically interpret \(T_{ab}\); indeed, the formulas for \(T_{ab}\) for fluids
and the fields considered above explicitly contain the metric. Thus, in general
relativity, one must solve simultaneously for the spacetime metric and the matter
distribution. This feature contributes to the difficulty of solving Einstein's
equation when sources are present.

The final remark concerns the equations of motion of matter. As we have presented
the theory, the equations of motion of particles, continuous matter, and fields are
postulated first, and then Einstein's equation relating the matter distribution to
the curvature of spacetime is given. However, Einstein's equation implies the
relation \(\nabla^aT_{ab}=0\), and this relation contains a great deal of
information on the behavior of matter. Indeed, for a perfect fluid, the relation
\(\nabla^aT_{ab}=0\) is the entire content of the equations of motion. Thus
for a fluid we may economize our assumptions by merely postulating the form of
\(T_{ab}\); the equations of motion of the fluid are already contained in
Einstein's equation. Notice that for a perfect fluid with \(P=0\), i.e., a fluid
composed of grains of \textquotedblleft dust\textquotedblright{} which exert no
forces upon each other, the fluid equation of motion~\eqref{eq:4.3.8} implied by
\(\nabla^aT_{ab}=0\) tells us that the individual dust particles move on
geodesics. More generally, it can be shown (Fock 1939; Geroch and Jang 1975)
that the relation \(\nabla^aT_{ab}=0\) implies that any sufficiently
\textquotedblleft small\textquotedblright{} body whose self-gravity is sufficiently
\textquotedblleft weak\textquotedblright{} must travel on a geodesic. Thus,
Einstein's equation alone actually implies the \emph{geodesic hypothesis} that the
world lines of test bodies are geodesics of the spacetime metric. This demonstrates
an important self-consistency of Einstein's equation with the basic framework of
general relativity. Note however, that bodies which are \textquotedblleft
large\textquotedblright{} enough to feel the tidal forces of the gravitational field
will deviate from geodesic motion. The equations of motion of such bodies also can
be found from the condition \(\nabla^aT_{ab}=0\) (Papapetrou 1951; Dixon
1974).
